\section{Метрические и нормированные пространства}
\label{sec:metric-and-normed-spaces}

\subsection{Метрические пространства}

\emph{Метрическим пространством} называется пара $(X,\rho)$, где $X$ ---
множество, а функция
\[
    \rho\colon X\times X\to[0,+\infty)
\]
удовлетворяет для любых $x,y,z\in X$ аксиомам:
\begin{enumerate}
    \item $\rho(x,y)=0$ тогда и только тогда, когда $x=y$;
    \item $\rho(x,y)=\rho(y,x)$;
    \item $\rho(x,z)\leq\rho(x,y)+\rho(y,z)$.
\end{enumerate}
Число $\rho(x,y)$ называется расстоянием между $x$ и $y$.

Открытый шар радиуса $r>0$ с центром $x$ определяется как
\[
    B(x,r)=\{y\in X:\rho(x,y)<r\}.
\]
Множество $G\subset X$ называется открытым, если с каждой своей точкой содержит
некоторый шар с центром в этой точке. Множество $F$ замкнуто, если его дополнение
открыто; эквивалентно, предел любой сходящейся последовательности элементов $F$
принадлежит $F$.

Примеры метрических пространств:
\[
    \mathbb R^n,\qquad
    \rho(x,y)=\sqrt{\sum_{k=1}^n(x_k-y_k)^2},
\]
и пространство непрерывных функций $C[a,b]$ с равномерной метрикой
\[
    \rho(f,g)=\max_{x\in[a,b]}|f(x)-g(x)|.
\]

\subsection{Сходимость, фундаментальные последовательности и полнота}

Последовательность $x_n$ сходится к $x\in X$, если
\[
    \rho(x_n,x)\to0.
\]
Последовательность называется \emph{фундаментальной}, или последовательностью
Коши, если
\[
    \forall\varepsilon>0\ \exists N:\quad
    \rho(x_n,x_m)<\varepsilon
    \quad\forall n,m>N.
\]
Любая сходящаяся последовательность фундаментальна, но обратное утверждение
верно не во всяком метрическом пространстве.

Метрическое пространство называется \emph{полным}, если всякая фундаментальная
последовательность в нём имеет предел, принадлежащий этому пространству.
Например, $\mathbb R^n$ полно.

Неполное метрическое пространство можно вложить как плотное подпространство в
полное пространство. Такое полное пространство называется его
\emph{пополнением}; оно определено с точностью до изометрии.

\subsection{Компактность}

Множество $K\subset X$ называется \emph{компактным}, если из любого его открытого
покрытия можно выделить конечное подпокрытие. В метрическом пространстве это
эквивалентно секвенциальной компактности:
\[
    \text{из любой последовательности }x_n\in K
    \text{ можно выделить подпоследовательность, сходящуюся к точке }K.
\]
Компактное множество в метрическом пространстве замкнуто и ограничено.
В $\mathbb R^n$ верно и обратное утверждение --- теорема Гейне--Бореля:
замкнутое и ограниченное множество компактно.

В бесконечномерных пространствах это уже неверно. Например, замкнутый единичный
шар в $\ell^2$ не компактен: последовательность стандартных базисных векторов
$e_n$ лежит в этом шаре, но
\[
    \|e_n-e_m\|_2=\sqrt2\qquad(n\neq m),
\]
поэтому из неё нельзя выделить сходящуюся подпоследовательность.

\subsection{Нормированные пространства}

Пусть $X$ --- линейное пространство. Функция $\|\cdot\|\colon X\to[0,+\infty)$
называется \emph{нормой}, если
\begin{enumerate}
    \item $\|\alpha x\|=|\alpha|\,\|x\|$;
    \item $\|x+y\|\leq\|x\|+\|y\|$;
    \item $\|x\|=0$ тогда и только тогда, когда $x=0$.
\end{enumerate}
Пара $(X,\|\cdot\|)$ называется нормированным пространством.

Любая норма порождает метрику
\[
    \rho(x,y)=\|x-y\|.
\]
Поэтому понятия сходимости, фундаментальности, открытости и полноты для
нормированного пространства определяются через эту метрику.

Примеры норм в $\mathbb R^n$:
\[
    \|x\|_1=\sum_{k=1}^n|x_k|,
    \qquad
    \|x\|_2=\left(\sum_{k=1}^n|x_k|^2\right)^{1/2},
    \qquad
    \|x\|_\infty=\max_k|x_k|.
\]
Для $C[a,b]$ естественна норма
\[
    \|f\|_\infty=\max_{x\in[a,b]}|f(x)|.
\]

\subsection{Банаховы пространства}

Полное нормированное пространство называется \emph{банаховым}. К банаховым
пространствам относятся
\[
    \mathbb R^n,\qquad \ell^p\ (1\leq p\leq\infty),\qquad
    C[a,b]\text{ с нормой }\|\cdot\|_\infty,
\]
а также пространства Лебега $L^p$ при $1\leq p\leq\infty$.

Любое конечномерное нормированное пространство полно. В бесконечномерном случае
полноту необходимо проверять отдельно; она является важным условием многих
теорем функционального анализа.

\subsection{Эквивалентность норм}

Две нормы $\|\cdot\|_1$ и $\|\cdot\|_2$ на одном линейном пространстве называются
\emph{эквивалентными}, если существуют константы $c,C>0$ такие, что
\[
    c\|x\|_1\leq\|x\|_2\leq C\|x\|_1,
    \qquad \forall x\in X.
\]
Эквивалентные нормы задают одну и ту же топологию: совпадают сходящиеся и
фундаментальные последовательности, открытые и замкнутые множества и свойство
полноты.

Важная теорема: \emph{на любом конечномерном линейном пространстве все нормы
эквивалентны}. В бесконечномерных пространствах различные нормы, вообще говоря,
могут быть неэквивалентны.


\subsection{Открытые и замкнутые множества, замыкание и граница}

В метрическом пространстве открытый шар имеет вид
\[
B_r(x)=\{y:\rho(x,y)<r\}.
\]
Множество $G$ открыто, если вместе с каждой точкой содержит некоторый шар вокруг
неё. Множество $F$ замкнуто, если содержит предел любой сходящейся последовательности
своих точек. В метрических пространствах эти определения эквивалентны тому, что
$X\setminus F$ открыто.

Замыкание $\overline A$ --- множество всех предельных точек и точек $A$;
внутренность $A^\circ$ --- наибольшее открытое множество, содержащееся в $A$;
граница
\[
    \partial A=\overline A\setminus A^\circ.
\]

\subsection{Компактность в метрических пространствах}

Для метрического пространства эквивалентны следующие формы компактности:
\begin{enumerate}
    \item из любого открытого покрытия можно выбрать конечное подпокрытие;
    \item из любой последовательности можно выделить сходящуюся подпоследовательность;
    \item пространство полно и вполне ограничено.
\end{enumerate}

В $\mathbb R^n$ действует теорема Гейне--Бореля: множество компактно тогда и
только тогда, когда оно замкнуто и ограничено. В бесконечномерном нормированном
пространстве это неверно: замкнутый единичный шар, вообще говоря, некомпактен.

\subsection{Конечномерные нормированные пространства}

На конечномерном линейном пространстве все нормы эквивалентны: для любых норм
$\|\cdot\|_a$ и $\|\cdot\|_b$ существуют $c,C>0$ такие, что
\[
    c\|x\|_a\le \|x\|_b\le C\|x\|_a.
\]
Следствия:
\begin{itemize}
    \item понятия сходимости и непрерывности не зависят от выбора нормы;
    \item любое конечномерное нормированное пространство полно;
    \item ограниченные последовательности имеют сходящиеся подпоследовательности
    после отождествления с $\mathbb R^n$ или $\mathbb C^n$.
\end{itemize}

\subsection{Пополнение пространства}

Всякое метрическое пространство можно изометрически вложить в полное метрическое
пространство, в котором его образ всюду плотен; такое пространство называется
пополнением. Для нормированного линейного пространства пополнение также можно
выбрать банаховым. Типичный пример: пространство непрерывных функций с некоторыми
нормами после пополнения приводит к пространствам $L^p$.

\subsection{Что важно сказать на экзамене}

Метрическое пространство задаётся расстоянием и тремя аксиомами метрики. Нужно
дать определения сходимости, фундаментальной последовательности и полноты.
Компактность в метрическом пространстве эквивалентна возможности выделить из
любой последовательности сходящуюся подпоследовательность. Норма на линейном
пространстве порождает метрику $\rho(x,y)=\|x-y\|$; полное нормированное
пространство называется банаховым. В конечномерном случае все нормы эквивалентны
и любое нормированное пространство полно.

\newpage
